% Nội dung 2.1--2.3: tiếng Việt, tiêu đề tiếng Anh.
\chapter{Literature Review and Background}
\label{chap:background}

\section{FIELD OVERVIEW}

\subsection{Technology Trends and Digital Transformation}

Một buổi tự học tiếng Nhật có thể gồm nhiều hoạt động: đọc bài, tra từ mới, nghe cách phát âm, làm bài tập và ôn lại kiến thức. Khi sử dụng các công cụ riêng cho từng việc, người học phải tự chuyển nội dung giữa chúng. Chẳng hạn, sau khi tra một từ trong bài đọc, họ cần ghi lại từ đó rồi tạo thẻ để ôn vào ngày hôm sau. Vấn đề đặt ra cho NihonGO là làm cho những bước này thuận tiện hơn và giúp người học tiếp tục sử dụng kiến thức vừa học.

Các công cụ học tập số cung cấp những cách hỗ trợ khác nhau. Từ điển giúp tìm nghĩa và cách dùng của từ. Flashcard, hay thẻ ghi nhớ, hỗ trợ ôn tập. Công cụ hội thoại tạo cơ hội thực hành nghe và trả lời. Trong đề tài này, nhóm xem xét ba hướng công nghệ liên quan trực tiếp đến những hoạt động đó.

\textbf{Thứ nhất, lưu kết quả để tổ chức ôn tập.} Với thẻ ghi nhớ điện tử, hệ thống có thể lưu người học đã nhớ hoặc quên từ nào. Anki sử dụng kết quả ôn để xác định thời điểm xuất hiện lại của thẻ; tài liệu chính thức giới thiệu các phương pháp lập lịch dựa trên SM-2 và FSRS~\cite{anki-scheduler}. Cách làm này là cơ sở tham khảo cho chức năng ôn tập của NihonGO. Thay vì yêu cầu người học tự tìm lại toàn bộ từ cũ, hệ thống cần giúp họ xác định những thẻ đến lúc ôn.

\textbf{Thứ hai, xử lý lời nói và âm thanh.} Tài liệu Gemini mô tả khả năng chuyển lời nói thành văn bản, phân tích âm thanh và trả lời câu hỏi về nội dung nghe được~\cite{gemini-audio}. Những khả năng này hỗ trợ xây dựng bài luyện nghe và chức năng trò chuyện. Ví dụ, trong tình huống gọi món, hệ thống cần nhận câu nói của người học, tạo câu trả lời của người phục vụ và phát lại bằng âm thanh. Độ chính xác của nội dung và thời gian chờ là hai yếu tố cần kiểm tra trong quá trình xây dựng.

\textbf{Thứ ba, sử dụng AI làm đối tác luyện hội thoại.} Duolingo đã giới thiệu Video Call với nhân vật Lily và công bố hỗ trợ tiếng Nhật~\cite{duolingo-videocall}. Điều này cho thấy hội thoại AI đã được đưa vào sản phẩm học ngoại ngữ. Đối với NihonGO, nhóm lựa chọn cách tổ chức hội thoại theo tình huống để người học biết mình cần làm gì trong cuộc trao đổi, chẳng hạn hỏi đường, mua hàng hoặc giới thiệu bản thân.

Từ ba hướng trên, NihonGO tập trung kết hợp bài luyện với công cụ tra cứu, ôn tập và xem lại kết quả. Người học cần biết bắt đầu ở đâu, hoàn thành bài như thế nào và nên ôn phần nào sau đó. Đây là những yêu cầu cụ thể để nhóm lựa chọn công nghệ và thiết kế chức năng.

\subsection{Related Market}

Theo khảo sát giáo dục tiếng Nhật ở nước ngoài năm 2024 của Japan Foundation, có 4.000.750 người học tại 143 quốc gia và vùng lãnh thổ. Việt Nam có 164.495 người học, đứng thứ sáu trong báo cáo. So với kỳ khảo sát năm 2021, số người học tại Việt Nam giảm 3,0\%~\cite{jpf2024}. Các số liệu này phản ánh số người học tại các cơ sở giảng dạy; hoạt động tự học cá nhân nằm ngoài phạm vi khảo sát.

Với NihonGO, người dùng được xác định là người học tiếng Nhật ở mọi trình độ, từ mới bắt đầu đến N1. Nhóm này gồm sinh viên, người đi làm và người học vì sở thích. Nhu cầu giữa các nhóm khác nhau: người mới cần làm quen chữ và mẫu câu; người đã có nền tảng cần ôn kiến thức, đọc hiểu và thực hành diễn đạt. Vì vậy, hệ thống cần phân loại nội dung theo trình độ để người học lựa chọn bài phù hợp.

Trong giai đoạn dùng thử, nhóm dự kiến tiếp cận sinh viên trước vì đây là đối tượng chiếm phần lớn mẫu khảo sát của đề tài. Sau đó, cần bổ sung người học trung cấp, cao cấp và người đi làm để đánh giá khả năng đáp ứng ở những nhóm còn lại. Phạm vi người dùng của sản phẩm vẫn bao gồm tất cả trình độ.

Về thương mại hóa, số người học trong báo cáo của Japan Foundation chỉ cung cấp bối cảnh, chưa cho biết có bao nhiêu người sẽ dùng hoặc trả phí cho NihonGO. Để xây dựng gói sử dụng, nhóm cần tìm hiểu thêm ngân sách của người học và chi phí vận hành từng chức năng. Nội dung này được tiếp tục phân tích trong Chương 3.

\section{MARKET DEMAND SURVEY AND RELATED PRODUCTS}

\subsection{Market Demand Survey}

\subsubsection{Survey Participants}

Khảo sát thu được 48 phản hồi từ ngày 16/09/2026 đến ngày 17/11/2026. Có 47 người là sinh viên và một người học vì sở thích. Về trình độ tự đánh giá, hai người vừa bắt đầu, 37 người ở mức N5--N4, sáu người ở mức N3 và ba người ở mức N2--N1~\cite{nihongo-survey2026}.

Bảng hỏi tìm hiểu ba nội dung chính: người học đang gặp khó khăn gì, họ đang sử dụng công cụ nào và họ muốn được hỗ trợ thêm điều gì. Các câu hỏi về khó khăn và công cụ cho phép chọn nhiều đáp án. Vì vậy, một người có thể được tính ở nhiều lựa chọn; tỷ lệ của mỗi lựa chọn được tính trên tổng số 48 người trả lời.

\subsubsection{Main Findings}

Kết quả khảo sát cho thấy nhu cầu cần chú ý trước hết là \textbf{ghi nhớ kiến thức}. Có 36 người chọn khó khăn về Kanji, tương đương 75,0\%, và 20 người chọn khó khăn về vốn từ nền tảng, tương đương 41,7\%. Một số ý kiến muốn học chữ qua bộ thủ, ví dụ và câu chuyện; các ý kiến khác muốn ôn tập ngắt quãng hoặc lưu lại từ đã tra. Những phản hồi này dẫn đến yêu cầu về tra cứu Kanji, lưu từ và ôn bằng flashcard trong NihonGO.

Nhu cầu thứ hai là \textbf{thực hành giao tiếp}. Một người trả lời cho biết đã học từ vựng và ngữ pháp nhưng vẫn khó tạo câu khi nói. Những người khác đề nghị hội thoại với AI, luyện phản xạ và nhận góp ý về cách dùng từ. Vì vậy, phần luyện nói của NihonGO cần tạo ra cuộc trao đổi hai chiều, trong đó người học tự đưa ra câu trả lời và nghe phản hồi theo tình huống.

Nhu cầu thứ ba là \textbf{hiểu cách dùng và biết chỗ cần sửa}. Các phản hồi đề cập đến thiếu giải thích ngữ cảnh và khó tự nhận biết lỗi phát âm. Từ đó, nhóm xác định phản hồi học tập cần chỉ rõ câu hoặc phần nội dung đang được nhận xét. Ví dụ, góp ý cho bài viết nên chỉ ra cách dùng từ chưa phù hợp và đề xuất cách diễn đạt khác. Riêng việc đánh giá phát âm cần kiểm tra khả năng xử lý âm thanh; bản chép lời chỉ cho biết hệ thống đã nhận dạng được nội dung nào.

Nhu cầu thứ tư là \textbf{duy trì việc học}. Có 35 người, tương đương 72,9\%, chọn khó khăn về động lực và tần suất. Các mong muốn được nêu gồm mục tiêu rõ ràng, hoạt động bớt nhàm chán và cảm giác hoàn thành sau khi học. Do đó, NihonGO cần hiển thị kết quả bài luyện và tiến độ để người học nhìn thấy những việc đã làm, đồng thời thuận tiện quay lại bài tiếp theo.

Đối với người học trình độ cao hơn, nhóm phân tích riêng chín phản hồi thuộc N3--N1 vì câu hỏi này dành cho đối tượng đó. Sáu người chọn khó khăn giao tiếp như người bản xứ, sáu người chọn tình trạng học chững lại và năm người chọn sắc thái ngữ pháp phức tạp. Kết quả gợi ý cần chuẩn bị thêm ngữ cảnh giao tiếp và giải thích cách diễn đạt cho người học trung cấp, cao cấp.

Ngoài nội dung học, một số phản hồi còn đề cập đến \textbf{chi phí}. Người trả lời mong muốn có chức năng giao tiếp miễn phí hoặc mức phí phù hợp với sinh viên. Đây là lý do để nhóm cân nhắc một gói cơ bản và giới hạn sử dụng AI rõ ràng. Mức giá cụ thể cần được khảo sát thêm vì bảng hỏi hiện tại chưa thu thập ngân sách của từng người.

Các số liệu và ý kiến trên được tổng hợp từ khảo sát của đề tài~\cite{nihongo-survey2026}. Do phần lớn người tham gia là sinh viên N5--N4, nhóm sử dụng kết quả để lựa chọn chức năng ban đầu và sẽ tiếp tục lấy ý kiến ở những trình độ khác. Hiệu quả của các chức năng được đánh giá khi người dùng trực tiếp sử dụng NihonGO.

\subsection{Related Products and Systems}

Nhóm lựa chọn Mazii, Anki, Quizlet, Duolingo và các công cụ trò chuyện AI để tìm hiểu vì chúng xuất hiện trong câu trả lời khảo sát. Thông tin chức năng dưới đây dựa trên tài liệu chính thức được tra cứu ngày 23/11/2026. Việc so sánh tập trung vào từng hoạt động học để xác định những điểm NihonGO cần tham khảo và hoàn thiện.

\subsubsection{Dictionary and Review Tools}

\textbf{Mazii} cung cấp tra từ, Kanji, ngữ pháp, tìm chữ viết tay và tra bằng camera. Nhà phát triển còn giới thiệu flashcard, đọc tin, luyện JLPT và các chức năng AI hỗ trợ dịch, hội thoại, phát âm, sửa ngữ pháp~\cite{mazii-app}. Mazii vì thế là sản phẩm cần đối chiếu ở cả tra cứu và luyện tập. Với NihonGO, nhóm chú ý đến cách người học tra một từ trong bài và lưu lại để ôn.

\textbf{Anki} được xem xét ở cách học bằng thẻ và lên lịch ôn. Phương pháp FSRS sử dụng lịch sử trả lời để xác định thời điểm ôn tiếp theo~\cite{anki-scheduler}. NihonGO có thể tham khảo cách ghi nhận mức nhớ của từng thẻ. Phần thực hành nghe, nói, đọc và viết cần tiếp tục tạo cơ hội cho người học dùng những từ đã ôn vào câu và tình huống.

\textbf{Quizlet} tổ chức nội dung thành các bộ thẻ gồm thuật ngữ và định nghĩa. Chế độ Flashcards hỗ trợ lật thẻ, phân loại nội dung đã biết, còn học và nghe âm thanh~\cite{quizlet-flashcards}. Chế độ Learn cung cấp các dạng câu hỏi theo mục tiêu và mức độ quen thuộc với bộ thẻ~\cite{quizlet-learn}. Nhóm tham khảo cách trình bày dễ thao tác và cách luyện cùng một bộ kiến thức bằng nhiều dạng bài.

\subsubsection{Conversation Tools}

\textbf{Duolingo} cung cấp Video Call với nhân vật AI Lily. Người học có thể trò chuyện, yêu cầu nhắc lại hoặc nói chậm hơn; chức năng được công bố có hỗ trợ tiếng Nhật~\cite{duolingo-videocall}. Đây là ví dụ để tham khảo khi thiết kế phần luyện hội thoại. Với NihonGO, nhóm cần kiểm tra người học có nghe được câu trả lời, hiểu tình huống và tiếp tục cuộc trao đổi hay không.

\textbf{Các công cụ trò chuyện AI} như ChatGPT, Gemini hoặc công cụ khác được 45/48 người trong khảo sát lựa chọn, chiếm 93,8\%~\cite{nihongo-survey2026}. Người học đã có thể sử dụng nhóm công cụ này cho những yêu cầu học tập của mình. Vì vậy, NihonGO cần làm rõ giá trị của việc học trong một ứng dụng riêng: có bài để chọn, có kết quả được lưu và có nội dung để tiếp tục ôn sau buổi học.

\begin{table}[htbp]
  \centering
  \caption{Những nội dung tham khảo từ các sản phẩm liên quan}
  \label{tab:related-product-tasks}
  \begin{tabular}{@{}p{2.1cm}p{5.6cm}p{6.7cm}@{}}
    \toprule
    \textbf{Sản phẩm} & \textbf{Hoạt động được tìm hiểu} & \textbf{Nội dung tham khảo cho NihonGO} \\
    \midrule
    Mazii & Tra cứu và luyện tập & Tra từ trong bài, lưu từ để ôn \\
    Anki & Ôn tập bằng thẻ ghi nhớ & Lưu mức nhớ và lên lịch ôn thẻ \\
    Quizlet & Học bằng bộ thẻ và câu hỏi & Giao diện ôn tập, các dạng bài \\
    Duolingo & Hội thoại với nhân vật AI & Cách bắt đầu và tiếp tục hội thoại \\
    Công cụ AI & Trao đổi theo yêu cầu của người học & Kết nối hỏi đáp với bài học và lịch sử \\
    \bottomrule
  \end{tabular}
  \par\smallskip
  {\fontsize{11pt}{14pt}\selectfont Nguồn: nhóm tổng hợp từ tài liệu sản phẩm và khảo sát đã dẫn tại Mục 1.2. Bảng trình bày nội dung tham khảo, không xếp hạng chất lượng sản phẩm.\par}
\end{table}

Qua việc tìm hiểu, nhóm xác định NihonGO cần tập trung vào quá trình học: chọn bài phù hợp, thực hành, xem phản hồi và ôn lại phần cần củng cố. Để đánh giá cách tổ chức này, nhóm sẽ kiểm tra người dùng có hoàn thành được bài và tìm lại kiến thức đã học hay không. Kết quả dùng thử sẽ cho biết những bước nào còn khó sử dụng và cần sửa.

\section{STRATEGIC ANALYSIS}

\subsection{SWOT Analysis}

Nhóm tổng hợp điểm mạnh, điểm yếu, cơ hội và thách thức của NihonGO trong Bảng~\ref{tab:swot-matrix}, dựa trên khảo sát nhu cầu, phạm vi sản phẩm và các công cụ học tiếng Nhật đã tìm hiểu.

\begin{table}[htbp]
\centering
\caption{Ma trận SWOT của NihonGO}
\label{tab:swot-matrix}
\begingroup
\setstretch{1}
\fontsize{11}{14}\selectfont
\renewcommand{\arraystretch}{1.1}
\begin{tabular}{|p{7.2cm}|p{7.2cm}|}
\hline
\begin{minipage}[t]{\linewidth}
\raggedright\smallskip
\textbf{S -- Strengths (Điểm mạnh)}
\begin{itemize}[leftmargin=1.2em,itemsep=4pt,parsep=0pt,topsep=5pt]
  \item \textbf{Xuất phát từ nhu cầu cụ thể:} khảo sát chỉ ra khó khăn về Kanji, giao tiếp và duy trì việc học, tạo căn cứ lựa chọn chức năng.
  \item \textbf{Dữ liệu học tập thống nhất:} tài khoản, bài luyện và lịch sử được tổ chức trong cùng hệ thống, hỗ trợ kết nối tra cứu, ôn tập và xem lại kết quả.
\end{itemize}\smallskip
\end{minipage}
&
\begin{minipage}[t]{\linewidth}
\raggedright\smallskip
\textbf{W -- Weaknesses (Điểm yếu)}
\begin{itemize}[leftmargin=1.2em,itemsep=4pt,parsep=0pt,topsep=5pt]
  \item \textbf{Ít dữ liệu nhu cầu ở trình độ cao:} chỉ 9/48 người khảo sát thuộc N3--N1; cần thêm phản hồi và kiểm tra học liệu theo trình độ.
  \item \textbf{Khối lượng phát triển lớn:} bốn kỹ năng đòi hỏi chuẩn bị nội dung, kiểm tra đáp án và phản hồi AI.
  \item \textbf{Thiếu đánh giá dài hạn:} chưa có kết quả theo dõi để xác định mức tiến bộ sau thời gian sử dụng.
\end{itemize}\smallskip
\end{minipage}
\\ \hline
\begin{minipage}[t]{\linewidth}
\raggedright\smallskip
\textbf{O -- Opportunities (Cơ hội)}
\begin{itemize}[leftmargin=1.2em,itemsep=4pt,parsep=0pt,topsep=5pt]
  \item \textbf{Nhu cầu phù hợp để dùng thử:} hội thoại AI, ôn tập ngắt quãng, giải thích theo ngữ cảnh và lưu kiến thức hay quên.
  \item \textbf{Có nền tảng AI hỗ trợ:} các dịch vụ như Gemini cung cấp khả năng xử lý âm thanh cho hoạt động luyện tập~\cite{gemini-audio}.
  \item \textbf{Kênh tiếp cận ban đầu:} sinh viên, nhóm tự học và câu lạc bộ tiếng Nhật có thể tham gia dùng thử, góp ý trực tiếp.
\end{itemize}\smallskip
\end{minipage}
&
\begin{minipage}[t]{\linewidth}
\raggedright\smallskip
\textbf{T -- Threats (Thách thức)}
\begin{itemize}[leftmargin=1.2em,itemsep=4pt,parsep=0pt,topsep=5pt]
  \item \textbf{Nhiều sản phẩm thay thế:} người học đã có Mazii, Anki, Quizlet, Duolingo và các công cụ AI.
  \item \textbf{Phụ thuộc dịch vụ bên ngoài:} giá, hạn mức hoặc gián đoạn API có thể ảnh hưởng chi phí và hoạt động học.
  \item \textbf{Rủi ro phản hồi sai:} lời giải thích chưa chính xác có thể khiến người học ghi nhớ sai; cần kiểm tra và cải thiện thường xuyên.
\end{itemize}\smallskip
\end{minipage}
\\ \hline
\end{tabular}
\endgroup
\end{table}

\Needspace{20cm}
\subsection{Porter's Five Forces Analysis}

Nhóm sử dụng mô hình năm áp lực cạnh tranh của Porter để xem xét môi trường phát triển NihonGO~\cite{hbs-fiveforces}. Bảng~\ref{tab:porter-five-forces} liên kết từng áp lực với nhận định chính và hướng xử lý của sản phẩm.

\begin{table}[htbp]
\centering
\caption{Năm áp lực cạnh tranh và hàm ý cho NihonGO}
\label{tab:porter-five-forces}
\begingroup
\setstretch{1}
\fontsize{11}{14}\selectfont
\renewcommand{\arraystretch}{1.15}
\begin{tabular}{@{}>{\raggedright\arraybackslash}p{4cm}>{\raggedright\arraybackslash}p{5.6cm}>{\raggedright\arraybackslash}p{5.4cm}@{}}
\toprule
\textbf{Áp lực cạnh tranh} & \textbf{Nhận định chính} & \textbf{Hàm ý cho NihonGO} \\
\midrule
\textbf{Đối thủ hiện tại}\newline Rivalry among Existing Competitors
& Người học đã có thói quen và dữ liệu trên Mazii, Anki, Quizlet hoặc Duolingo. Tính năng mới cần đem lại lợi ích rõ ràng trong một hoạt động học cụ thể.
& \textbf{Ưu tiên trải nghiệm học hoàn chỉnh:} tra và lưu từ thuận tiện, hiểu góp ý sau bài làm, tìm lại kết quả để ôn tập. \\
\midrule
\textbf{Đối thủ mới}\newline Threat of New Entrants
& Dịch vụ AI sẵn có giúp các nhóm khác xây dựng ứng dụng hội thoại. Chỉ tích hợp AI vào giao diện khó tạo lợi thế lâu dài.
& \textbf{Đầu tư vào học liệu và chất lượng bài luyện:} kiểm tra nội dung, phân loại trình độ và cải thiện từ phản hồi thực tế. \\
\midrule
\textbf{Sản phẩm thay thế}\newline Threat of Substitutes
& Giáo trình, video, giáo viên và câu lạc bộ cùng đáp ứng nhu cầu học tiếng Nhật. Người học có thể kết hợp nhiều hình thức.
& \textbf{Làm rõ vai trò hỗ trợ tự luyện:} kết nối tra cứu, thực hành và ôn tập trong những buổi học phù hợp với thời gian sẵn có. \\
\midrule
\textbf{Nhà cung cấp}\newline Bargaining Power of Suppliers
& Giá, hạn mức và gián đoạn dịch vụ AI, xử lý tiếng nói hoặc hạ tầng ảnh hưởng đến chi phí và khả năng vận hành. Học liệu cần có nguồn và điều kiện sử dụng rõ ràng.
& \textbf{Giảm phụ thuộc và kiểm soát chi phí:} tách kết nối dịch vụ khỏi nghiệp vụ, quản lý hạn mức, xử lý gián đoạn và kiểm tra nguồn học liệu. \\
\midrule
\textbf{Khách hàng}\newline Bargaining Power of Users
& Người học có nhiều công cụ để lựa chọn; khảo sát ghi nhận nhu cầu về mức phí phù hợp với sinh viên.
& \textbf{Ưu tiên giá trị và chi phí phù hợp:} làm rõ chức năng, hạn mức từng gói; khảo sát ngân sách và đo chi phí phục vụ trước khi chốt giá. \\
\bottomrule
\end{tabular}
\endgroup
\end{table}

\noindent\textbf{Hàm ý chiến lược.}
Nhóm tập trung vào ba ưu tiên: hoàn thiện trải nghiệm học cốt lõi, xây dựng học liệu phù hợp trình độ và kiểm soát chi phí AI. Các ưu tiên này được cụ thể hóa trong yêu cầu hệ thống tại Mục~\ref{sec:system-requirements} và định hướng mô hình kinh doanh ở Chương 4.

\section{SYSTEM REQUIREMENTS}
\label{sec:system-requirements}
\subsection{Functional Requirements}
\label{sec:functional-requirements}
% Phạm vi được đối chiếu với mã nguồn trên develop, cập nhật ngày 25/11/2026.
% Nguồn: 02_FunctionalRequirement.xlsx, sheet Project Requirements Template.
% Mã yêu cầu chính thức của báo cáo: FR-01, FR-02, ...; mỗi dòng một mã.
% Các mã UC trong ghi chú nội bộ bên dưới chỉ tham chiếu tài liệu đặc tả gốc.
%
% Các phần chưa đưa vào báo cáo:
% - UC-S19--S32: phòng hội thoại tự do/ghép cặp người thật; không đánh đồng
%   với phòng game LiveKit. Migration hoặc model riêng lẻ không đủ làm bằng chứng.
% - UC-W16: chưa có API/giao diện xóa bài đã nộp.
% - UC-V22: chưa thấy game nối từ; UC-V26--V28: chưa thấy xếp hạng liên kỳ/huy hiệu.
% - UC-R6: chưa có bản dịch toàn bài; UC-R14, UC-R16: chưa có luồng danh sách
%   lịch sử đọc/đánh dấu hoàn thành riêng. UC-R15 chỉ khôi phục kết quả gần nhất.
% - UC-CM07--CM08, UC-CM12--CM19: chưa thấy đủ luồng đổi mật khẩu khi đã đăng nhập,
%   xóa tài khoản, nhắc học chung, hộp thông báo, thanh toán/quản lý gói trả phí.
%   Nhắc ôn flashcard native được giữ riêng ở UC-V20.
% Các mô tả thu hẹp đáng chú ý:
% - UC-CM01/03: chỉ ghi email/mật khẩu, không cam kết mọi nhà cung cấp OAuth.
% - UC-CM06: chỉ xem thông tin; UC-CM10: mục tiêu và phút/ngày, không ghi hạn thi.
% - UC-V1--V2: từ vựng theo bài và tra Kanji; chưa có kho ngữ pháp độc lập.
% - UC-V8--V12: yêu thích từ, tạo/sửa bộ; không ghi yêu thích Kanji/ngữ pháp,
%   mô tả bộ hay xóa bộ vì API chưa hỗ trợ đủ.
% - UC-S11--S12: nhận xét transcript, không chấm âm vị/phát âm/trọng âm.
% - UC-S16: chỉ mở lại phiên còn hoạt động, không có thao tác pause riêng.
% - UC-W7: tự lưu web cần mạng, không bảo đảm khôi phục phần chưa lưu khi offline.
% - UC-W19--W21: so sánh hai bài cùng đề, câu thêm/bớt; không ghi đếm lỗi
%   hay đánh dấu thay đổi theo từng ký tự giữa các phiên bản.
%
% Bằng chứng chính: backend/{SPEAKING,WRITING,READING,FLASHCARDS,SHADOWING,
% VOCABULARY,DICTIONARY,GAMES,LEARNING}.md; routes.py tương ứng;
% frontend/src/screens/{WritingWorkspaceScreen,ReadingScreen,ShadowingScreen,
% AccountScreen,GameLibraryScreen,SpeakingScenariosScreen}.tsx;
% frontend/src/components/{WritingVersionComparison,DictionaryWordActions}.tsx;
% frontend/src/hooks/useWritingDocuments.ts.

Các yêu cầu chức năng dưới đây được tổng hợp từ bảng đặc tả của đề tài~\cite{nihongo-functional2026} và giới hạn theo các thao tác đã có trong phiên bản được trình bày. Mỗi yêu cầu mô tả việc người học có thể thực hiện và kết quả hệ thống cần cung cấp. Mỗi yêu cầu được gán một mã duy nhất theo định dạng FR-xx (Functional Requirement -- yêu cầu chức năng), đánh số liên tục từ FR-01 theo thứ tự trình bày trong bảng. Các mã này được sử dụng thống nhất khi tham chiếu yêu cầu chức năng trong báo cáo.

\begingroup
\renewcommand{\arraystretch}{1.15}
\setlength{\LTleft}{0pt}
\setlength{\LTright}{\fill}
\begin{longtable}{@{}>{\raggedright\arraybackslash}p{3.4cm}p{12cm}@{}}
  \caption{Các yêu cầu chức năng trong phạm vi hệ thống NihonGO}
  \label{tab:functional-requirements}\\
  \toprule
  \textbf{Mã yêu cầu} & \textbf{Yêu cầu chức năng} \\
  \midrule
  \endfirsthead

  \multicolumn{2}{c}{\tablename~\thetable{} (tiếp theo)}\\
  \toprule
  \textbf{Mã yêu cầu} & \textbf{Yêu cầu chức năng} \\
  \midrule
  \endhead

  \midrule
  \multicolumn{2}{r}{\textit{Tiếp ở trang sau}}\\
  \endfoot

  \bottomrule
  \endlastfoot

  \multicolumn{2}{@{}l}{\textbf{1. Tài khoản và thiết lập học tập}}\\*
  FR-01 & \textbf{Đăng ký và xác thực.} Người dùng nhập thông tin đăng ký và mật khẩu. Hệ thống gửi mã xác thực qua email, kiểm tra mã và tạo tài khoản khi thông tin hợp lệ. \\
  FR-02 & \textbf{Đăng nhập và đăng xuất.} Người dùng đăng nhập bằng email và mật khẩu; đăng xuất khi kết thúc sử dụng. \\
  FR-03 & \textbf{Khôi phục mật khẩu.} Người dùng yêu cầu mã xác nhận và đặt mật khẩu mới. Hệ thống kiểm tra mã trước khi cập nhật mật khẩu. \\
  FR-04 & \textbf{Xem thông tin tài khoản.} Người dùng xem họ tên, thông tin liên hệ và trạng thái xác thực của tài khoản. \\
  FR-05 & \textbf{Thiết lập việc học.} Người dùng tự chọn trình độ từ mới bắt đầu đến N1, mục tiêu và số phút học mỗi ngày. Hệ thống lưu các lựa chọn và cho phép cập nhật sau đó. \\
  \midrule

  \multicolumn{2}{@{}l}{\textbf{2. Học từ vựng và tra cứu}}\\*
  FR-06 & \textbf{Chọn bài học.} Người dùng chọn trình độ N5--N1, xem danh sách bài từ vựng và mở bài muốn học. \\
  FR-07 & \textbf{Xem từ và nghe cách đọc.} Hệ thống hiển thị từ, cách đọc và nghĩa; người dùng mở thông tin bổ sung khi có và nghe cách đọc bằng chức năng phát âm. \\
  FR-08 & \textbf{Tra cứu bằng từ khóa.} Người dùng nhập chữ Nhật, kana, romaji hoặc nghĩa tiếng Việt để tìm từ; chọn kết quả để xem chi tiết. Với Kanji, hệ thống hiển thị thông tin như âm đọc, nghĩa và thứ tự nét khi có trong dữ liệu. \\
  FR-09 & \textbf{Tra chữ viết tay.} Người dùng vẽ một ký tự tiếng Nhật. Hệ thống đưa ra các ký tự gợi ý để người dùng chọn và tra cứu. \\
  FR-10 & \textbf{Xem từ đã tra.} Hệ thống lưu các mục từ và Kanji đã mở trong lịch sử gần đây của tài khoản. Người dùng chọn một mục để xem lại. \\
  FR-11 & \textbf{Quản lý từ yêu thích.} Người dùng thêm hoặc bỏ một từ khỏi danh sách yêu thích và mở danh sách đó để ôn tập. \\
  FR-12 & \textbf{Tạo và sửa bộ thẻ.} Người dùng đặt tên bộ thẻ, nhập từ và nghĩa, chỉnh sửa tên hoặc thêm, sửa, bỏ thẻ trong bộ. Từ tra trong từ điển có thể được thêm vào bộ tự tạo. \\
  \midrule

  \multicolumn{2}{@{}l}{\textbf{3. Flashcard và lịch ôn tập}}\\*
  FR-13 & \textbf{Học bằng flashcard.} Người dùng chọn bộ có sẵn, bộ tự tạo hoặc từ yêu thích; lật thẻ và chọn một trong bốn mức nhớ. Khi kết thúc, hệ thống hiển thị số thẻ đã học và kết quả tự đánh giá. \\
  FR-14 & \textbf{Ôn thẻ đến hạn.} Hệ thống tính lần ôn tiếp theo từ mức nhớ đã chọn, hiển thị số thẻ cần ôn theo ngày và cho phép bắt đầu phiên ôn các thẻ đã đến hạn. \\
  FR-15 & \textbf{Nhắc ôn trên điện thoại.} Người dùng bật nhắc ôn và chọn giờ. Khi được cấp quyền, ứng dụng hẹn thông báo cục bộ theo lịch thẻ. Chức năng nhắc này áp dụng cho ứng dụng di động. \\
  FR-16 & \textbf{Xem tiến độ ôn tập.} Người dùng xem số từ đã tự đánh giá, các từ còn yếu và số thẻ cần ôn để lựa chọn hoạt động tiếp theo. \\
  \midrule

  \multicolumn{2}{@{}l}{\textbf{4. Nghe và Shadowing}}\\*
  FR-17 & \textbf{Chọn nội dung nghe.} Người dùng chọn video trong thư viện hoặc nhập liên kết YouTube được hỗ trợ. Hệ thống hiển thị trạng thái chuẩn bị và mở bài khi dữ liệu phụ đề đã sẵn sàng. \\
  FR-18 & \textbf{Luyện theo đoạn.} Người dùng nghe liên tục hoặc chọn từng đoạn để nghe và nhắc lại; có thể phát lại, lặp đoạn và thay đổi tốc độ phát. \\
  FR-19 & \textbf{Xem phụ đề và bản dịch.} Phụ đề tiếng Nhật hiển thị theo thời gian video. Người dùng có thể ẩn, hiện phụ đề và bật bản dịch tiếng Việt; hệ thống thông báo trạng thái khi bản dịch đang được chuẩn bị. \\
  FR-20 & \textbf{Tra từ trong phụ đề.} Người dùng chọn từ để xem cách đọc, nghĩa và sử dụng các thao tác lưu từ của công cụ từ điển. \\
  FR-21 & \textbf{Nghe điền từ.} Người dùng nghe đoạn video, nhập phần còn thiếu trong câu và kiểm tra đúng, sai theo nội dung phụ đề. \\
  FR-22 & \textbf{Trả lời câu hỏi nghe hiểu.} Người dùng chọn trình độ và mở bộ 5 câu hỏi theo phụ đề; hệ thống dùng bộ đã lưu hoặc yêu cầu AI tạo khi chưa có. Hệ thống dừng video tại mốc câu hỏi, hiển thị kết quả và giải thích sau khi trả lời. Kết quả trả lời được xem lại trong phiên đang mở. \\
  \midrule

  \multicolumn{2}{@{}l}{\textbf{5. Hội thoại với AI}}\\*
  FR-23 & \textbf{Chọn tình huống.} Người dùng xem, lọc tình huống theo trình độ và chủ đề; đọc bối cảnh, vai, mục tiêu và bắt đầu phiên hội thoại. \\
  FR-24 & \textbf{Trao đổi với AI.} Người dùng nói hoặc nhập văn bản và chuyển cách nhập trong phiên. AI trả lời theo vai và bối cảnh của tình huống, kèm âm thanh phản hồi. \\
  FR-25 & \textbf{Lưu nội dung hội thoại.} Hệ thống lưu từng lượt trao đổi và thời gian. Người dùng lựa chọn lưu âm thanh để có thể nghe lại khi bản ghi đã được lưu thành công. \\
  FR-26 & \textbf{Nhận xét sau hội thoại.} Với phiên có đủ dữ liệu, hệ thống nhận xét ngữ pháp, từ vựng và độ tự nhiên dựa trên bản chép lời; hiển thị điểm và gợi ý diễn đạt. Nhận xét này không chấm phát âm. \\
  FR-27 & \textbf{Quản lý lịch sử hội thoại.} Người dùng xem danh sách phiên, mở transcript và nhận xét, nghe bản ghi có sẵn hoặc xóa phiên của mình. \\
  FR-28 & \textbf{Tiếp tục hoặc kết thúc.} Người dùng mở lại phiên còn hoạt động để tiếp tục, kết thúc phiên khi luyện xong hoặc tạo phiên mới với tình huống đã chọn. \\
  \midrule

  \multicolumn{2}{@{}l}{\textbf{6. Đọc hiểu}}\\*
  FR-29 & \textbf{Chọn bài đọc.} Người dùng xem danh sách, lọc theo trình độ và dạng bài, tìm bằng từ khóa và mở nội dung bài đọc. \\
  FR-30 & \textbf{Hỗ trợ đọc từ.} Người dùng bật, tắt furigana và chọn từ trong bài để xem cách đọc, nghĩa qua cửa sổ tra cứu. \\
  FR-31 & \textbf{Xem hướng dẫn đọc.} Hệ thống hiển thị các mẹo đọc hiểu phù hợp với dấu hiệu trong bài để người học tham khảo. \\
  FR-32 & \textbf{Làm và xem lại câu hỏi.} Người dùng chọn đáp án và nộp bài. Hệ thống tính điểm, hiển thị đáp án, cho phép yêu cầu AI giải thích và khôi phục kết quả gần nhất khi mở lại đúng phiên bản bài đọc. \\
  FR-33 & \textbf{Phân tích bài đọc.} Theo yêu cầu, AI cung cấp tóm tắt, từ vựng và ngữ pháp đáng chú ý. Người dùng chọn một câu để xem phân tích cấu trúc và phần nội dung đang được giải thích. \\
  FR-34 & \textbf{Lưu từ khi đọc.} Người dùng tra một từ trong bài và lưu vào danh sách yêu thích hoặc bộ thẻ cá nhân qua cửa sổ từ điển. \\
  \midrule

  \multicolumn{2}{@{}l}{\textbf{7. Luyện viết}}\\*
  FR-35 & \textbf{Chọn đề viết.} Người dùng xem và lọc danh sách đề, đọc yêu cầu, xem gợi ý và mở vùng soạn bài. Bộ lọc có các mức N5--N1; bài được hiển thị theo nội dung hiện có. \\
  FR-36 & \textbf{Soạn và quản lý bản nháp.} Người dùng nhập bài, theo dõi số ký tự, lưu, mở lại hoặc xóa nháp. Bản web tự lưu khi ngừng nhập và có kết nối; hệ thống báo nếu lưu không thành công. \\
  FR-37 & \textbf{Nộp bài và xem góp ý.} Hệ thống kiểm tra nội dung trước khi nộp, yêu cầu AI nhận xét và hiển thị điểm theo tiêu chí. Người dùng xem bản gốc, các đoạn được góp ý, lý do sửa và bản sửa đề xuất. \\
  FR-38 & \textbf{Xem và làm lại bài.} Người dùng mở bài đã nộp để xem nội dung và nhận xét, sau đó chỉnh sửa để nộp lần mới. Mỗi lần nộp được lưu riêng, giữ nguyên bài đã nộp trước đó. \\
  FR-39 & \textbf{So sánh hai lần viết.} Người dùng chọn hai bài cùng đề để đối chiếu nội dung, điểm và nhận xét. Hệ thống liệt kê các câu được thêm hoặc bỏ giữa hai bài. \\
  \midrule

  \multicolumn{2}{@{}l}{\textbf{8. Trò chơi từ vựng}}\\*
  FR-40 & \textbf{Chọn trò chơi và bộ từ.} Người dùng chọn trò trong kho game, trình độ và nguồn từ. Với chế độ nhiều người, người dùng tạo phòng hoặc tham gia bằng mã phòng. \\
  FR-41 & \textbf{Luyện từ qua trò chơi.} Người dùng trả lời câu hỏi trắc nghiệm trong Công thành Momo, nhập cách đọc trong Đối kháng từ vựng hoặc giành lượt đọc trong Giành mic khi điều kiện âm thanh được đáp ứng. \\
  FR-42 & \textbf{Xem kết quả trận.} Hệ thống kiểm tra câu trả lời theo luật của trò chơi, tính điểm và hiển thị kết quả khi trận kết thúc. \\
\end{longtable}
\endgroup

\subsection{Non-Functional Requirements}
\section{THEORETICAL BACKGROUND}
\section{CHAPTER SUMMARY}
